
\begin{proposition}[Proved global conductor Euler bound]\label{cdp:st:C.2.VII.2}

Let $S$ be a compact irreducible reduced hypersurface in
a smooth complex threefold, with $\omega_S$ torsion and
$\chi(S,\OO_S)=0$. Let $\eta:T\to S$ be normalization,
let $\Gamma$ be its conductor image scheme, and put
$C=\Gamma_{\mathrm{red}}$. Then
\[
 \chi(C,\OO_C)\ge-1.
\tag{C.VII.3}\label{cdp:eq:C.VII.3}
\]
Suppose in addition that the minimal resolution
$\pi:Z\to T$ is a class VII surface. Then
\[
 e(C)\ge0.
\tag{C.VII.4}\label{cdp:eq:C.VII.4}
\]
The conductor scheme may be nonreduced and its generic
transverse exponent is unrestricted.
\end{proposition}
\begin{proof}
If the conductor is empty, both conclusions are immediate.
We therefore assume that it is nonempty.
Put $L=\omega_S$ and
$\varepsilon=h^0(S,L)\in\{0,1\}$.
A nonzero section of a torsion line on a compact
connected reduced space trivializes it, by taking a
torsion-trivializing power. Serre duality and
$\chi(\OO_S)=0$ therefore give
$h^1(S,\OO_S)=1+\varepsilon$, with
$\varepsilon=1$ exactly when $L=\OO_S$.
The normalization quotient
$\mathcal Q=\eta_*\OO_T/\OO_S$ has
$H^0(\mathcal Q)\hookrightarrow H^1(S,\OO_S)$,
because $S$ and $T$ have the same constant holomorphic
functions. The actual conductor duality of \ref{cdp:st:A.5.12} is
\[
 \mathcal Q=\omega_\Gamma\otimes L^{-1}|_\Gamma .
\]
Serre duality on the Cohen--Macaulay conductor curve
gives $h^1(\Gamma,L|_\Gamma)\le1+\varepsilon$.
The restriction onto $C$ is a quotient of coherent
sheaves on a one-dimensional space. Its kernel has
no second cohomology, so
$h^1(C,L|_C)\le1+\varepsilon$.
Since $L|_C$ is torsion, it has degree zero on every
irreducible component and
$\chi(C,L|_C)=\chi(C,\OO_C)$.
If $\varepsilon=0$, Riemann--Roch gives
$\chi(\OO_C)=h^0(C,L|_C)-h^1(C,L|_C)\ge-1$.
If $\varepsilon=1$, then $L=\OO_S$ and
$h^0(C,\OO_C)\ge1$, while $h^1(C,\OO_C)\le2$.
This proves \eqref{cdp:eq:C.VII.3}, including all nilpotents of
the original conductor scheme in the comparison.

Now assume that $Z$ is class VII. \ref{cdp:st:A.5.3} gives a
nonzero effective integral divisor $D$ with
$K_Z+D=(\eta\pi)^*\omega_S$.
Each curve in the reduced strict conductor on $Z$
is rational or elliptic by \ref{cdp:st:C.2.VII.1}. Its finite image
in $C$ has normalization of genus at most one, by
Riemann--Hurwitz. Let $n$ be the number of irreducible
components of $C$, let $g_i\in\{0,1\}$ be the genera
of their normalizations, and let $r_x$ and $\delta_x$
be the branch number and delta invariant at a singular
point of $C$. The exact normalization sequence and
the topological normalization map give
\[
 \chi(\OO_C)=\sum_i(1-g_i)-\sum_x\delta_x,\qquad
 e(C)=\sum_i(2-2g_i)-\sum_x(r_x-1).
\tag{C.VII.5}\label{cdp:eq:C.VII.5}
\]
For every reduced curve germ,
$\delta_x\ge r_x-1$: its normalization permits
independent constants on the branches, whereas a
function on the original germ has one common constant.
If $n_0$ is the number of rational-normalization
components, \eqref{cdp:eq:C.VII.3}--\eqref{cdp:eq:C.VII.5} imply
$e(C)\ge n_0-1$.

The only possible negative value is consequently
$e(C)=-1$, with every normalized component elliptic.
Then \eqref{cdp:eq:C.VII.3} says $\sum\delta_x\le1$.
Negativity requires two branches, so there is exactly
one singular point, it has two branches and
$\delta=1$, and the local curve ring is
\[
 \{(a(u),b(v))\in\C\{u\}\oplus\C\{v\}:a(0)=b(0)\}.
\tag{C.VII.6}\label{cdp:eq:C.VII.6}
\]
Indeed that subring already has colength one in the
normalization, and the curve ring is contained in it
with the same colength. Thus the point is an ordinary
node. Its two branches come either from distinct
elliptic normalizations or from two distinct points
of one elliptic normalization.

We show that this node is incompatible with the
specified hypersurface and resolution maps. At least
one normalization point $t\in T$ over the node receives
preimages of both local branches of $C$.
Otherwise partition the irreducible local branches
of the hypersurface germ $S$ according to the one
conductor branch their normalization receives. Both
parts are nonempty. Two hypersurface branches from
opposite parts intersect in a curve: in the regular
three-dimensional local ambient ring, the ideal of
their two distinct equations has height two.
That intersection curve is contained in the conductor
image and therefore is one of its two branches.
Both surface branches must receive it, contrary to
the partition. With only one local surface branch
the assertion is immediate.

Every $\pi$-exceptional curve not contained in
$D_{\mathrm{red}}$ is disjoint from it. To see this,
write $E$ for such a curve. Minimality of the
resolution and adjunction give $K_Z\cdot E\ge0$,
while $D\cdot E\ge0$, and
$(K_Z+D)\cdot E=0$ forces equality. In particular
$D\cdot E=0$. A resolution fibre over a point of
the conductor must meet its strict transform, by
properness. Since its fibre is connected, this
disjointness proves
$\pi^{-1}(t)\subset D_{\mathrm{red}}$ for the
normalization point just found.

The strict curves over the two elliptic-normalized
branches are elliptic: a rational curve cannot map
nonconstantly onto an elliptic curve. They meet the
connected fibre $\pi^{-1}(t)$. If the curves are
distinct, one connected component of $D_{\mathrm{red}}$
therefore contains two elliptic curves. If they are
the same elliptic curve, their meeting points with
the fibre are distinct. This follows from the exact
factorization of a map from a normal curve through
the normalization of its image: one point cannot
map to both points over the node in \eqref{cdp:eq:C.VII.6}.
The connected resolution fibre consequently gives
a path in $D_{\mathrm{red}}$ joining these distinct
points. Both possibilities contradict C.2.VII.1:
the first has two positive-genus components, and
the second adds a graph cycle to an elliptic
component, giving arithmetic genus at least two.
This excludes the only negative possibility and
proves \eqref{cdp:eq:C.VII.4}.
\end{proof}


